\chapter{System Design and Architecture}\doublespacing

\section{Overview}
System design transforms the abstract functional and non-functional requirements defined in the SRS into a cohesive, robust, and maintainable software architecture. The design of the \textbf{LightAngo} compiler follows modern compiler engineering principles, emphasizing high cohesion, loose coupling, modular phase separation, and clear interface abstractions.

This chapter details the high-level architecture, module decomposition, use case interactions, data flow models (Level 0 and Level 1 DFDs), execution flowcharts, and activity representations that govern the internal working of the LightAngo compiler.

\section{Use Case Diagram}
The primary actors interacting with the LightAngo system are the \textbf{Developer (Programmer)}, the \textbf{Compiler Front-End Engine}, and the \textbf{Underlying Operating System / Toolchain} (NASM Assembler and GNU Linker).

Figure~\ref{fig:usecase} illustrates the principal use cases and operational interactions.

\begin{figure}[h!]
\centering
\begin{picture}(400,240)
% Outer boundary
\put(80,10){\framebox(280,220){}}
\put(90,215){\textbf{LightAngo Compiler System Boundary}}

% User Actor
\put(20,130){\circle{16}} % head
\put(20,122){\line(0,-1){24}} % body
\put(6,110){\line(1,0){28}} % arms
\put(20,98){\line(-1,-1){18}} % left leg
\put(20,98){\line(1,-1){18}} % right leg
\put(0,70){\textbf{Developer}}

% OS Actor
\put(380,130){\circle{16}} % head
\put(380,122){\line(0,-1){24}} % body
\put(366,110){\line(1,0){28}} % arms
\put(380,98){\line(-1,-1){18}} % left leg
\put(380,98){\line(1,-1){18}} % right leg
\put(365,70){\textbf{OS/Toolchain}}

% Use cases (ovals)
\put(110,180){\oval(140,24)} \put(115,177){\small UC-01: Write Source Code}
\put(110,145){\oval(140,24)} \put(115,142){\small UC-02: Submit File to CLI}
\put(230,175){\oval(150,22)} \put(160,172){\small UC-03: Lexical Tokenization}
\put(230,145){\oval(150,22)} \put(160,142){\small UC-04: Recursive AST Parsing}
\put(230,115){\oval(150,22)} \put(160,112){\small UC-05: Comment Validation}
\put(230,85){\oval(150,22)}  \put(160,82){\small UC-06: Semantic Scope Check}
\put(230,55){\oval(150,22)}  \put(160,52){\small UC-07: x86-64 NASM Gen}
\put(230,25){\oval(150,22)}  \put(160,22){\small UC-08: Assemble \& Link}

% Connection lines
\put(25,120){\vector(1,1){45}}
\put(25,115){\vector(1,0){45}}
\put(305,30){\vector(1,1){55}}

\end{picture}
\caption{Use Case Diagram of the LightAngo Compiler System} \label{fig:usecase}
\end{figure}

The developer provides input source files through the CLI. The compiler verifies syntactic structure, checks documentation comments, performs type and scope resolution, and upon success, orchestrates NASM and `ld` to output a native Linux binary.

\section{Data Flow Diagrams (DFD)}
Data Flow Diagrams visually represent how source representations and semantic data transform as they traverse the compiler subsystems.

\subsection{Level 0 Data Flow Diagram (Context Diagram)}
The Level 0 DFD (Figure~\ref{fig:dfd0}) models the entire compilation system as a single cohesive process receiving raw source input and delivering either formatted diagnostic errors or a compiled executable machine binary.

\begin{figure}[h!]
\centering
\begin{picture}(420,130)
% User entity
\put(10,50){\framebox(80,40){\shortstack{\textbf{User} \\ \small (Programmer)}}}

% Main Process
\put(170,40){\oval(120,60)}
\put(185,68){\textbf{LightAngo}}
\put(178,54){\textbf{Compilation}}
\put(195,40){\textbf{Engine}}

% Output Entity
\put(340,50){\framebox(70,40){\shortstack{\textbf{Target} \\ \textbf{Binary}}}}

% Data flows
\put(90,75){\vector(1,0){80}} \put(95,82){\small Source Code (.lango)}
\put(230,70){\vector(1,0){110}} \put(235,77){\small Native ELF Binary}
\put(170,50){\vector(-1,0){80}} \put(95,38){\small Error Diagnostics}

\end{picture}
\caption{Level 0 Data Flow Diagram (Context Level)} \label{fig:dfd0}
\end{figure}

\subsection{Level 1 Data Flow Diagram}
The Level 1 DFD (Figure~\ref{fig:dfd1}) decomposes the compiler into its internal phase-wise transformations and data stores (Symbol Table and Comment Registry).

\begin{figure}[h!]
\centering
\begin{picture}(430,280)
% Modules
\put(20,240){\framebox(90,25){\small 1.0 Source Reader}}
\put(170,240){\framebox(90,25){\small 2.0 Lexical Scanner}}
\put(170,180){\framebox(90,25){\small 3.0 AST Parser}}
\put(170,120){\framebox(90,25){\small 4.0 Semantic Check}}
\put(170,60){\framebox(90,25){\small 5.0 NASM Generator}}
\put(170,5){\framebox(90,25){\small 6.0 Assembler/Linker}}

% Stores
\put(310,180){\framebox(110,25){\small D1: Comment Registry}}
\put(310,120){\framebox(110,25){\small D2: Scoped Symbol Table}}

% Diagnostics Output
\put(20,90){\framebox(90,40){\shortstack{\small Diagnostic \\ \small Engine}}}

% Arrows
\put(110,252){\vector(1,0){60}} \put(115,257){\tiny Characters}
\put(215,240){\vector(0,-1){35}} \put(220,225){\tiny Token Stream}
\put(260,245){\vector(1,-1){50}} \put(280,240){\tiny Comments}
\put(310,192){\vector(-1,0){50}} \put(265,197){\tiny Doc Rules}
\put(215,180){\vector(0,-1){35}} \put(220,165){\tiny AST Nodes}
\put(260,132){\vector(1,0){50}}
\put(310,132){\vector(-1,0){50}} \put(265,137){\tiny Scope Lookup}
\put(215,120){\vector(0,-1){35}} \put(220,105){\tiny Validated AST}
\put(215,60){\vector(0,-1){30}} \put(220,45){\tiny x86-64 Assembly}

\put(170,190){\vector(-1,-1){60}} \put(115,160){\tiny Syntax Err}
\put(170,130){\vector(-1,0){60}} \put(115,135){\tiny Semantic Err}
\end{picture}
\caption{Level 1 Data Flow Diagram of the LightAngo Compiler Pipeline} \label{fig:dfd1}
\end{figure}

\section{System Flowchart}
The compilation execution flowchart (Figure~\ref{fig:flowchart}) traces the algorithmic decision tree from file input to terminal binary execution.

\begin{figure}[h!]
\centering
\begin{picture}(380,300)
% Start
\put(190,285){\oval(60,20)} \put(177,282){\small Start}
\put(190,275){\vector(0,-1){20}}

% Read file
\put(145,235){\framebox(90,20){\small Read Source File}}
\put(190,235){\vector(0,-1){20}}

% Lexer
\put(145,195){\framebox(90,20){\small Tokenize Stream}}
\put(190,195){\vector(0,-1){20}}

% Decision: Loop / Function?
\put(190,155){\oval(110,20)} \put(140,152){\small Has Loop / Function?}
\put(190,145){\vector(0,-1){20}} \put(195,135){\tiny YES}
\put(245,155){\vector(1,0){30}} \put(255,160){\tiny NO}

% Comment Check Decision
\put(190,105){\oval(100,20)} \put(145,102){\small Valid Comment?}
\put(190,95){\vector(0,-1){20}} \put(195,85){\tiny YES}
\put(140,105){\vector(-1,0){40}} \put(105,110){\tiny NO}

% Error Box
\put(40,95){\framebox(60,20){\small Emit Error}}
\put(70,95){\vector(0,-1){75}}

% Code gen
\put(275,105){\framebox(90,20){\small Parse \& Semantic}}
\put(320,105){\vector(0,-1){30}}
\put(275,55){\framebox(90,20){\small Emit x86 Assembly}}
\put(320,55){\vector(0,-1){35}}

% Assemble
\put(145,15){\framebox(90,20){\small Assemble \& Link}}
\put(275,25){\vector(-1,0){40}}
\put(145,25){\vector(-1,0){45}}

% Exit
\put(70,10){\oval(60,18)} \put(58,7){\small Exit}

\end{picture}
\caption{System Flowchart of the Compilation and Verification Flow} \label{fig:flowchart}
\end{figure}

\section{Activity Diagram}
The activity diagram captures the parallel validation paths executed during compilation: lexical verification and documentation linting proceed concurrently prior to semantic graph construction and intermediate representation synthesis.

\section{Layered System Architecture}
The LightAngo compiler architecture is partitioned into three decoupled functional layers:
\begin{enumerate}
    \item \textbf{Input / Presentation Layer:} Manages CLI arguments, file buffer reads, source coordinate mapping, and formatted console diagnostic rendering.
    \item \textbf{Core Processing Layer:} Comprises the analytical front-end (Lexer, Recursive-Descent Parser, AST Builder, Comment Verification Engine) and semantic middle-end (Scoped Symbol Table, Type Checker, Constant Folding Optimizer).
    \item \textbf{Synthesis / Target Layer:} Comprises the AST Assembly Emitter, x86-64 stack frame manager, NASM assembler interface, and GNU linker execution.
\end{enumerate}

\section{System Module Decomposition}
Table~\ref{tab:system_modules} summarizes the architectural modules, their underlying methodology, and algorithmic implementation within the LightAngo codebase.

\begin{table}[h!]
\centering
\caption{System Module Specifications} \label{tab:system_modules}
\setstretch{1.3}
\begin{tabular}{|p{3.2cm}|p{3.8cm}|p{7.5cm}|}
\hline
\textbf{Module Name} & \textbf{Method / Technique} & \textbf{Algorithmic Responsibility} \\
\hline
\textbf{Lexical Scanner} & Deterministic Finite Automata (DFA) State Machine & Scans character stream, tokenizes keywords, operators, literals, and captures comment metadata. \\ \hline
\textbf{AST Parser} & Recursive Descent with Operator Precedence & Constructs Abstract Syntax Tree nodes for statements, binary expressions, and control flow blocks. \\ \hline
\textbf{Comment Validator} & Syntax-Bound Lookback Verification & Gating rule engine verifying non-empty comments precede loop and subroutine declarations. \\ \hline
\textbf{Semantic Analyzer} & Scoped Symbol Table Hierarchy & Type verification, identifier binding, scope boundary validation, and undeclared variable detection. \\ \hline
\textbf{Optimizer} & Tree-Walk Constant Folding & Evaluates static compile-time arithmetic expressions and prunes dead branches. \\ \hline
\textbf{Target Code Generator} & Stack-Based x86-64 Assembly Synthesis & Emits native NASM instructions (`push`, `pop`, `mov`, `cmp`, `jmp`, `syscall`) with stack alignment. \\ \hline
\textbf{Toolchain Driver} & System V ABI \& ELF Linker Subsystem & Invokes NASM assembler (`-f elf64`) and GNU linker (`ld`) to produce standalone native binaries. \\ \hline
\textbf{Error Diagnostics} & Source-Mapped Error Formatter & Prints color-coded, line/column-precise diagnostic feedback and remediation hints. \\ \hline
\end{tabular}
\end{table}

\newpage
